% 2011 年上海春季卷第 11 题：程序框图（i=0,s=66；s>0 时反复 s=s-10、i=i+1）
% 来源 PDF：春季高考/2011/2011春季上海.pdf，第 1 页（原刊 2-32 页）第 2 栏。
% 重绘依据：原图为「开始 → i=0,s=66 → 判断 s≤0?」，否分支向右执行 s=s-10、i=i+1 后
%   回到判断入口，是分支向下「输出 i」并「结束」。
\begin{tikzpicture}[scale=0.9, semithick, >=Stealth,
  block/.style={draw, minimum width=2.6cm, minimum height=0.7cm, align=center},
  dec/.style={draw, diamond, aspect=3.0, inner sep=1pt, align=center},
  arr/.style={->}]
  \node (start) [draw, rounded corners=6pt, minimum width=1.5cm, minimum height=0.55cm] {开始};
  \node (init)  [block, below=0.5cm of start] {$i=0,\ s=66$};
  \node (dec)   [dec, below=0.9cm of init] {$s\leqslant0$?};
  \node (out)   [draw, trapezium, trapezium left angle=70, trapezium right angle=110,
                 minimum width=2.4cm, minimum height=0.65cm, below=0.9cm of dec] {输出 $i$};
  \node (end)   [draw, rounded corners=6pt, minimum width=1.5cm, minimum height=0.55cm,
                 below=0.5cm of out] {结束};
  \node (sub)   [block, right=1.6cm of dec] {$s=s-10$};
  \node (add)   [block, above=0.7cm of sub] {$i=i+1$};
  \coordinate (join) at ($(init)!0.5!(dec)$);
  \draw[arr] (start) -- (init);
  \draw[arr] (init) -- (dec);
  \draw[arr] (dec) -- node[right] {是} (out);
  \draw[arr] (out) -- (end);
  \draw[arr] (dec.east) -- node[above] {否} (sub.west);
  \draw[arr] (sub.north) -- (add.south);
  \draw (add.west) -- (join -| add.west) -- (join);
  \draw[arr] (join -| add.west) -- (join);
\end{tikzpicture}
